\section{Introduction}

Designing a fault-tolerant quantum architecture requires choosing an error-correcting code, a syndrome-extraction procedure, and protocols for implementing logical gates. Comparing these choices often requires circuit-level simulation under physically motivated noise models, since their logical error rates and thresholds are generally difficult to determine analytically. Efficient Clifford simulation tools, such as \texttt{stim} \cite{gidney2021stim}, have made it possible to evaluate and compare a wide range of quantum memory experiments and logical Clifford operations.

Extending this capability to non-Clifford operations, which are required for universal quantum computation, remains a challenge. Numerous fault-tolerant protocols have been proposed, either to implement such operations directly or to prepare encoded \emph{magic states} that enable them through state injection. These include magic state distillation \cite{bravyi2005universal,bravyi2012magic,campbell2017unified,litinski2019magic}, cultivation \cite{gidney2024magic,vaknin2026high}, code switching \cite{beverland2021cost,bombin2014gauge,kubica2015universal,tan2025singleshot}, and pieceable fault tolerance \cite{yoder2016universal}. Recent constructions have strengthened the prospects for preparing magic states at high rates, including constant-overhead distillation protocols based on codes with nonvanishing encoding rate \cite{golowich2024asymptotically, golowich2024quantum, he2025asymptotically, li2026transversal, nguyen2024good, scruby2025quantum, steinberg2025universal, tan2025singleshot, wills2024constantoverhead, wills2026improved, xu2025batched, zhu2025nonclifford}. Such approaches reduce the overhead per output by processing many magic states within a growing code block, making their performance at large numbers of logical qubits particularly important to assess.

The need to evaluate these protocols has motivated recent simulation algorithms and software libraries tailored to non-Clifford quantum error correction \cite{haenel2026tsim,chase2026clifft,umbrarescu2026syqma,fang2026symft,surti2026efficient}. These methods substantially extend the range of accessible circuits, but generally retain potentially exponential runtime dependencies. In particular, Surti, Daguerre, and Kim \cite{surti2026efficient} establish polynomial scaling in the number of physical qubits for broad families of magic state preparation protocols when the number $k$ of logical qubits is fixed. Their complexity also depends on a stabilizer or Pauli decomposition of the target logical state, whose size can grow exponentially with $k$. Consequently, efficient simulation at fixed $k$ does not guarantee efficient simulation of high-rate protocols, where the number of logical outputs grows with the block size.

This motivates simulation methods that scale polynomially in both the physical and logical system sizes, even as the number of non-Clifford operations grows. In this work, we identify a broad class of circuits relevant to fault-tolerant magic state preparation that admit such simulation under stochastic Pauli noise.
\subsection{Summary of results}

We identify a class of non-Clifford quantum circuits that admit classical simulation in time polynomial in the number of qubits and the circuit size. The supported operations are:
\begin{itemize}
    \item Preparation of $|0\rangle$ and $|+\rangle$ states;
    \item Pauli, $\CNOT$, and $\T$ gates;
    \item Arbitrary $Z$-type Pauli measurements and compatible $X$-type Pauli measurements.
\end{itemize}
To describe the measurement restriction, we first identify the states generated by the unitary operations above. A \emph{third-order phase polynomial (3PP) state} is a uniform-magnitude superposition over an affine subspace of computational-basis labels (that is, the solution set of a system of linear equations over $\F_2$), with phases specified by a third-order phase polynomial over $\mathbb{Z}_8$. These are exactly the states obtainable from the specified preparations using Pauli, $\CNOT$, and $\T$ gates. A projective Pauli measurement is \emph{compatible} with a 3PP state if the post-measurement state is a 3PP state for every outcome having nonzero probability. We call circuits composed of the operations above, with each measurement compatible on every reachable branch, \emph{3PP circuits}.

The simulation maintains an exact, efficient representation of the conditional state throughout the circuit and samples intermediate measurement outcomes according to their probabilities. It also supports exact evaluation of computational-basis amplitudes in a fixed global-phase convention, exact Pauli expectation values, computational-basis sampling, and equality testing between represented states up to global phase. Stochastic Pauli noise is incorporated by sampling faults and evolving the corresponding conditional states, provided that measurements remain compatible along the noisy trajectories.

Alongside the phase-polynomial description, we develop an equivalent \emph{diagonal-Clifford-and-Pauli (DCP) stabilizer formalism}, in which states are uniquely specified, up to global phase, by stabilizing operators generated by Pauli gates and diagonal Clifford gates. We prove that a state is a DCP stabilizer state if and only if it is a 3PP state. We establish structural properties of DCP stabilizer groups, construct efficient tableau representations, and develop efficient update rules for both the tableau and phase-polynomial descriptions.

We give necessary and sufficient conditions for an $X$-type Pauli measurement to be compatible, together with polynomial-time compatibility tests and conditional-state updates. These results extend efficient simulation beyond unitary CNOT-dihedral circuits---circuits made of $\T$, $\CNOT$, and Pauli gates \cite{cross2016scalable}---to circuits with intermediate measurements and polynomial-time classical feedforward.

Our main result for quantum error correction reduces compatibility under noise to a condition on the ideal circuit.
\Cref{thm:iterative-measurements} shows that, if every ideal $X$ measurement is deterministic or anticommutes with a current $Z$-type stabilizer, then every measurement remains compatible under arbitrary inserted DCP faults, provided that outcome-dependent operations are restricted to DCP corrections.
Thus, these ideal-circuit conditions certify exact polynomial-time simulation of the entire noisy protocol, without a separate analysis of noisy branches.
This covers broad families of magic state preparation protocols, including distillation, cultivation, and code switching, as well as non-Clifford logical operations interleaved with syndrome extraction.

Beyond circuits that preserve the DCP state class, we introduce \emph{DCP rank}: the minimum number of DCP stabilizer states needed to express a state as a linear combination. Such decompositions extend the framework to universal circuit simulation, with a cost that depends on the number of terms. Since every Pauli stabilizer state is a DCP state, DCP rank is no greater than stabilizer rank, although it can still grow exponentially.

Finally, we provide \texttt{merlin}~\cite{koch2026merlin}, an open-source implementation of our simulation methods. We compare its runtime and memory usage with existing non-Clifford simulation libraries on distillation, cultivation, and code switching circuits.
The benchmarks demonstrate improved scaling as the number of logical outputs grows in Bravyi-Haah distillation \cite{bravyi2012magic}, and simulation of a recent code switching circuit based on a $[[81,3,5]]$ code \cite{li2026transversal} on which all other tested libraries exhaust memory or fail.
They also illustrate the complementary strengths of existing methods, which achieve lower per-shot runtimes on the single-output distillation and cultivation examples.

\subsection{Related work}

\subsubsection{Non-Clifford simulation algorithms}

Classical simulation beyond Clifford circuits includes stabilizer decompositions \cite{bravyi2016improved,bravyi2019simulation}, quasiprobability and Pauli-propagation methods \cite{pashayan2015estimating,rall2019simulation}, and hybrid stabilizer--tensor-network representations \cite{masot-llima2024stabilizer,nakhl2025stabilizer}.
Recent work has adapted and optimized these approaches for quantum error correction.
Tsim \cite{haenel2026tsim} combines stabilizer decompositions with ZX-calculus simplification and parametric compilation \cite{kissinger2022classical,sutcliffe2025fast}, while cultivation-specific decompositions demonstrate that circuit structure can substantially reduce the number of required Clifford terms \cite{wan2025cutting,wan2026simulating}.
SOFT \cite{li2025soft}, Clifft \cite{chase2026clifft}, and SymFT \cite{fang2026symft} combine stabilizer or Clifford-frame descriptions with explicit coefficients for non-stabilizer degrees of freedom.
Other approaches include sparse Pauli-frame representations with truncation \cite{tuloup2026computing} and polynomial-memory symbolic simulation through auxiliary qubits, as implemented in SyQMA \cite{umbrarescu2026syqma}.

More specifically exploiting magic-state preparation, Surti, Daguerre, and Kim \cite{surti2026efficient} propagate circuit-level Pauli faults into Clifford errors, with complexity depending on a decomposition of the target magic state.
Magdalena de la Fuente et al.~\cite{fuente2026highthreshold} develop a complementary spacetime approach for efficiently sampling syndrome statistics in a class of non-Clifford CSS-type circuits.
Our approach establishes polynomial-time simulation for circuits preserving the DCP state class, including compatible intermediate measurements, without requiring a small $\T$ count or a small number of logical magic states.

\subsubsection{Simulating \texorpdfstring{$\CNOT+\T$}{CNOT+T} circuits}

The unitary part of 3PP circuits consists of $\CNOT$, $\T$, and Pauli gates acting on qubits initialized in $\ket{0}$ or $\ket{+}$. Circuits with this gate set have been studied in several contexts in the quantum information literature.

This gate set is closely related to the Z-network formulation of Instantaneous Quantum Polynomial-time (IQP) circuits \cite[Section 5]{shepherd2009instantaneous}. With $\ket{+}$ inputs and $X$-basis measurements, Z-networks generated by $\CNOT$, $X$, and $e^{i\theta Z}$ are equivalent to IQP sampling up to efficient classical post-processing. Restricting the rotation angles to integer multiples of $\pi/8$ gives our gate set, up to global phase. For this choice of angles, Shepherd subsequently showed that marginal distributions on logarithmically many output bits can be computed efficiently \cite[Theorem 4]{shepherd2010binary}.

The unitary group generated by $\CNOT$, $\T$, and Pauli gates was subsequently named the CNOT-dihedral group and studied in the contexts of randomized benchmarking and circuit synthesis \cite{cross2016scalable, amy2018finite, garion2020synthesis}. Efficient polynomial representations and composition algorithms were established in \cite{cross2016scalable}, while \cite{amy2018finite} developed a finite presentation using phase polynomials.

Here, we characterize when $X$ measurements preserve the class of 3PP states and derive efficient conditional-state update rules, enabling simulation of circuits with intermediate measurements satisfying these conditions.

\subsubsection{Generalized stabilizer formalisms}

Our DCP stabilizer formalism can be seen as a special case of the \textit{monomial stabilizer formalism} introduced by Van den Nest \cite{nest2011monomial}, whose stabilizers are products of diagonal unitaries and permutation matrices. DCP operators have this form because Pauli $X$ operators are permutation matrices. In this framework, each stabilizer state is an equal-magnitude superposition of computational-basis states, connected by the stabilizers' permutation actions. For DCP states, these actions are binary translations, so the support is an affine subspace. Our phase-polynomial representation satisfies the monomial framework's sufficient conditions for efficient computational-basis sampling and approximate local expectation values. We additionally provide exact Pauli expectation values and efficient gate and compatible-measurement updates.

More directly, the XS formalism \cite{ni2015noncommuting} uses tensor products of $X$ and $\S$ operators with eighth-root global phases. XS states already admit affine-support, third-order phase-polynomial descriptions, but with additional constraints linking the polynomials to the support. By allowing $\CZ$ factors in stabilizers, DCP removes these constraints and captures all third-order phase-polynomial states.

The XP formalism \cite{webster2022xp} further generalizes XS to arbitrary phase precision. Its measurement analysis \cite[Section 7 and Appendix F]{webster2022xp} provides outcome probabilities for XP measurements and updates for diagonal Pauli measurements. For Pauli $X$ measurements, its relative-phase calculation corresponds to the discrete derivative of the phase polynomial in our formulation. However, it does not classify when non-diagonal measurements preserve the XP state class. We provide such a classification for DCP states, together with polynomial-time compatibility tests and post-measurement updates.

\vspace{\baselineskip}
\paragraph*{Note added.}
While completing this manuscript, we became aware of two independent related works on efficient simulation of noisy non-Clifford quantum error correction.
Takada, Bartlett, and Williamson \cite{takada2026exact} develop a Clifford-stabilizer method closely related to our tableau approach.
Here, we also describe the states using phase polynomials, establish the equivalence of the two representations, and classify the measurements preserving this state class.
Magdalena de la Fuente and Scruby \cite{fuente2026efficiently} instead use a spacetime framework to sample measurement outcomes and bound logical failure rates without explicitly tracking the logical state, whereas our method maintains the conditional state throughout the circuit.
Additionally, we provide an open-source implementation and extend our framework to universal simulation through DCP state decompositions (\Cref{sec:universal-simulation}).
Exploring the precise relationships and practical tradeoffs between these approaches is an interesting direction for future work.

\subsection{Reading guide}

We develop the simulation framework through two complementary descriptions of the same states and circuits.
We begin with third-order phase-polynomial states and compatible circuits in \Cref{sec:phase-polynomial-states}, using this representation to construct a simulator in \Cref{sec:simulator-phase-polynomial-representation}.
We then recast these states in the DCP stabilizer formalism in \Cref{sec:dcp-stabilizer-formalism}, leading to an alternative, tableau-based simulator in \Cref{sec:simulator-dcp-representation}.
With both descriptions in place, \Cref{sec:noisy-simulations} establishes their applicability to noisy quantum error-correction circuits, and \Cref{sec:benchmarking-results} benchmarks our implementation, \texttt{merlin}, against existing simulators.
Finally, \Cref{sec:universal-simulation} extends the approach beyond the efficiently simulable class through DCP decompositions, before we discuss further directions in \Cref{sec:discussion}.
Technical proofs and algorithmic details are collected in the appendices.

Readers primarily interested in understanding the simulator can focus on \Cref{sec:phase-polynomial-states,sec:simulator-phase-polynomial-representation} and defer the DCP formalism on a first reading. The stabilizer picture in \Cref{sec:dcp-stabilizer-formalism,sec:simulator-dcp-representation} provides further structural insights and an alternative way to understand the same states and circuits, as well as an alternative implementation of the simulator.
